\documentclass[landscape]{article}

% Page layout
\usepackage[margin=1in]{geometry}   % 1-inch margins
\usepackage{graphicx}
\usepackage{subcaption}
\usepackage{pgffor} % Required for the \foreach loop


\setlength{\parindent}{0pt}         % No paragraph indentation
\setlength{\parskip}{0.25em}        % Space between paragraphs
\raggedbottom   

% Define a safe macro to handle the logic for each patient
\long\def\processpatient#1{%
    \ifnum#1=3\relax
        % Skip 3
    \else\ifnum#1=111\relax
        % Skip 111
    \else\ifnum#1=178\relax
        % Skip 178
    \else\ifnum#1=228\relax
        % Skip 228
    \else
        % If it is none of the numbers above, compile the page normally
        \begin{figure}[p] 
            \centering
            \includegraphics[width=\linewidth, height=12cm, keepaspectratio]{Patient_#1}
            \caption{Patient #1}
            \label{fig:patient-#1}
        \end{figure}
        \clearpage
    \fi\fi\fi\fi % You need one \fi for every \ifnum opened
}

\begin{document}

% Define the path to your images folders
\graphicspath{{/Users/sillouz2/Documents/MATLAB/RunModel 5/PatientFiguresOpt}}

\title{Patient Plots Catalog}
\author{Logan L, Charlotte B, Simon I, Noah R, Morgan S}
\date{\today}
\maketitle
\maketitle % Generates the title on the first page

% ==========================================
% 1. FIRST PATIENT (Placed on the title page)
% ==========================================
\begin{figure}[h!] 
    \centering
    \includegraphics[width=\linewidth, height=10cm, keepaspectratio]{Patient_1}
    \caption{Patient 1}
    \label{fig:patient-1}
\end{figure}
\clearpage % Starts a fresh page ONLY AFTER the first patient is printed

\foreach \j in {1,...,343} {
    \processpatient{\j}
}

\end{document}